% =====================================================================
%  Solar Filament Segmentation (IEEE BigData Cup 2026) -- team report
%  Build:  pdflatex main.tex ; bibtex main ; pdflatex main.tex ; pdflatex main.tex
%  Length budget: 4 pages (ACM acmart / sigconf).
% =====================================================================
\documentclass[sigconf]{acmart}

% preamble.tex -- shared packages and macros for the Solar Filament report.
% Included by main.tex; kept separate so that the document body can be edited
% without touching the package setup.

\usepackage{booktabs}
\usepackage{multirow}
\usepackage{amsmath}
\usepackage{graphicx}
\usepackage{seqsplit}

% This TeX Live installation has no libertine / newtxmath / inconsolata, so
% acmart falls back to non-scalable Type 3 fonts in the front-matter blocks.
% microtype's automatic font expansion then fails on those blocks, so the
% expansion is switched off (the fallback fonts are used for the whole
% document anyway).
\microtypesetup{expansion=false}

% acmart puts \shorttitle in the running head but this class version does not
% define the macro, so it is defined here (no argument) to keep the header
% from overflowing into the conference name.
\def\shorttitle{Solar Filament Segmentation}

% Small helper: a number printed with a fixed number of decimals.
\newcommand{\num}[2]{\ensuremath{#1#2}}

% Compact "delta" column with an explicit sign.
\newcommand{\dlt}[1]{\ensuremath{\Delta\mathrm{PQ} = #1}}

% Keep floats inside their column where possible and avoid float-induced
% overflow of the two-column layout.
\setlength{\textfloatsep}{10pt plus 2pt minus 2pt}
\setlength{\floatsep}{8pt plus 2pt minus 2pt}
\setlength{\abovecaptionskip}{4pt}
\setlength{\belowcaptionskip}{2pt}

% Give the justifier a little extra room, so that long numbers and inline
% formulas do not push lines out of the column.
\setlength{\emergencystretch}{1.5em}

% Breakable typewriter text for artifact file names (they are long and must
% not overflow the two-column measure).
\newcommand{\fn}[1]{\texttt{\seqsplit{#1}}}

% Data-availability note used in the reproducibility paragraph.
\newcommand{\artifacts}{All validation numbers quoted here were produced by the
retained evaluation artifacts of the project (per-variant evaluation JSONs,
the seed table, the post-processing and ensemble sweep tables, and the
captured public-leaderboard listing); each claim in this paper is traceable to
one of those files.}


\acmConference[Solar Filament Segmentation Challenge 2026]{IEEE BigData Cup 2026}{November 2026}{Online}
\settopmatter{printacmref=false} % libertine/newtxmath are not installed in this TeX Live,
                                 % so the ACM reference-format block cannot be typeset
\renewcommand\footnotetextcopyrightpermission[1]{} % no copyright block in this template

\title{Consensus Pseudo-Labels for Solar\\
Filament Instance Segmentation}

\author{Team feasylol}

\keywords{instance segmentation, solar filaments, consensus pseudo-labels,
U-Net, panoptic quality, sunspot and filament analysis}
\ccsdesc[500]{Computing methodologies~Image processing}
\ccsdesc[300]{Applied computing~Astronomy}

\begin{document}

\begin{abstract}
Solar filaments are thin, low-contrast dark structures in full-disk solar
images; the challenge asks for them as \emph{instances}, not as a foreground
mask. The organiser's metric is a row-wise panoptic-quality (PQ) score: every
annotator--image row of the ground truth is scored against the same predicted
instances, so a prediction must agree with the annotators simultaneously,
rather than with one of them. We first attacked the task with a
hand-crafted consensus detector, then replaced it by learning. Our pipeline
turns the public multi-annotator polygons into per-pixel consensus
pseudo-labels (a support count plus a per-class vote), trains a
7.76M-parameter U-Net on 1024$\times$1024 maps, extracts instances by
thresholding, a $7\times7$ morphological closing, 4-connectivity labelling and
an area floor with a per-image relaxation fallback, and averages four seeds
with horizontal-flip test-time augmentation. On a
fixed 88-image validation split (220 annotator--image rows, 1{,}514
ground-truth masks, 21 annotators) the local replication of the official score
rises from 0.2201 to 0.3449, and the public-leaderboard score from 0.18 to
0.29. We also report negative results (class splitting, geometric watershed, a
wider network) and the optimism ($\approx 0.015$ PQ) of using one shared
validation split for monitoring, thresholding and seed selection.
\end{abstract}

\maketitle

\section{Introduction}
Filaments are magnetic features that appear as dark, elongated structures on
the solar disk. Segmenting them is hard for two independent reasons: they are
weakly contrasted against the photosphere, and different human annotators
trace the same filament differently. The Solar Filament Segmentation Challenge
of the IEEE BigData Cup 2026 provides 1{,}154 annotator--image annotation
entries over 707 distinct full-disk images (8{,}199 single-polygon instances,
four category labels: Left, Right, Unidentifiable and Ambiguous), while the 180
test images ship without public annotations. Annotator disagreement is
therefore part of the task definition, not noise around it.

The metric makes this explicit. The organiser scores each
\emph{annotator--image} row separately and pools the rows: a prediction is a
true positive if it overlaps a ground-truth instance of that row with IoU
strictly above 0.5, matched pairs are counted (not unique matches), and the
score is $\sum\mathrm{IoU}(\mathrm{TP})/(|TP|+0.5|FP|+0.5|FN|)$~\cite{notebook}.
A prediction that satisfies only one
annotator is a true positive in that row and a false positive in every other
row, so the metric rewards shapes that several annotators agree on --- the
design principle of our second system.

Our first system was a hand-crafted detector that fused several pseudo-annotator
scales into an instance-anchored consensus and thresholded a local-contrast
image; it reached a local PQ of 0.2201 and a public score of 0.18. Its ceiling
was obvious: every improvement came from re-tuning hand-made thresholds. All
comparisons below are between our own variants, measured with one protocol on
one split; we make no claim about methods we did not run.

We contribute: (i) a consensus pseudo-label construction that is derived from
the public polygons only, together with the measurement that motivates it
(Section~\ref{sec:consensus}); (ii) a U-Net recipe and an instance-extraction
pipeline with a relaxation fallback (Section~\ref{sec:method}); (iii) an
evaluation protocol that keeps the validation images out of training and
quantifies the selection optimism that this shared split introduces
(Section~\ref{sec:eval}); and
(iv) an ablation table that reports the variants that failed
(Section~\ref{sec:ablation}).

\section{Methodology}
\label{sec:method}

\subsection{Consensus pseudo-labels}
\label{sec:consensus}
The public training polygons are the only supervision available. We rasterised
every annotator's polygons onto the native $2048\times2048$ canvas with the
official COCO routine (polygon-to-mask followed by mask merging) and built
three per-pixel maps per image:
\emph{support} (how many annotators' instances cover the pixel, $0..M$ where
$M$ is the number of annotators of that image), \emph{votes} (four channels,
one per category, counting the annotators covering the pixel with that
category) and \emph{class} (the arg-max of the votes, ties broken in the fixed
order Left $>$ Right $>$ Unidentifiable $>$ Ambiguous). The maps are built at
the native $2048\times2048$ resolution and, for the $1024\times1024$ models,
area-averaged in $k\times k$ blocks with $k=2048/1024=2$; the soft target is
that average over its per-image maximum, the class map its arg-max over the
four channels. The image is resized $2048^2\to1024^2$ bilinearly, so one
example is a whole disk, not a sub-crop. The label files also carry a
$512\times512$ view ($4\times4$ block mean of support, block sum of votes)
that the earlier $512$-pixel model consumed and the $1024$-pixel training code
never reads. The category statistics of this
release are Left 2{,}535, Right 2{,}590, Unidentifiable 3{,}074 and Ambiguous
0 instances; the Ambiguous channel is kept but is empty, and 3{,}459{,}165
pixels are covered by two or more annotators.

Why a consensus instead of a single annotator? The annotations genuinely
disagree: over 3{,}152 greedily matched instance pairs the median IoU is
0.5762 and only 68\% of the pairs exceed IoU 0.5. We therefore measured
detectors that copy one annotator against detectors that vote, grouping the
rows by the number of annotators per image. On two-annotator images, copying an
annotator gives PQ 0.6802 and 0.6794, whereas the support $\geq 2$ vote gives
0.7090. On three-annotator images the copy baseline averages 0.5556 and the
same vote reaches 0.6248. Leave-one-out over the 151 three-annotator images
confirms the direction (0.3462 for the union rule against 0.3334 for copying
one annotator). The consensus is therefore not a smoothing heuristic: scored
against all annotators at once it is closer to them collectively than any
single annotation, which is what the row-wise metric rewards.

\subsection{Network and training}
We train a five-level U-Net with 32 base channels (7{,}759{,}653 parameters) on
the $1024\times1024$ consensus maps, with random flips and intensity jitter.
The optimiser is Adam with a
learning rate of $10^{-3}$ on a cosine schedule, mixed precision, batch size 4
and 50 epochs; the loss is a soft binary cross-entropy on the support map plus
$0.5\times$ cross-entropy on the class map. A base-32 run takes 28.95--30.16 minutes
on a single Tesla V100S (155 optimiser steps per epoch). We ran 15 such runs
over 13 distinct seeds (the two diverged seeds 7 and 15 were re-run); 6 of
them diverged with a non-finite loss and are reported as such rather than
silently dropped.

\subsection{Instance extraction}
At inference the $1024\times1024$ probability map is bilinearly upsampled to
the native $2048\times2048$ canvas, on which every step below acts: the map is
thresholded at 0.70, closed with a
$7\times7$ morphological kernel, labelled with 4-connectivity, and components
smaller than 600 pixels are discarded. Because an empty prediction for an image
is catastrophic under a row-wise metric, an image that would end up empty is
re-processed along a relaxation chain (threshold 0.55 with area floor 200,
then threshold 0.45, then the largest component alone), so that no test image
is left without an instance. On the submitted ensemble, 5 of the 180 test
images needed that fallback and the produced file passed the submission
contract check with zero issues.

\subsection{Ensembling}
Four independently trained seeds (three with the same recipe but different
seeds, plus the seed used for the single-model submission) are averaged in
probability space, each with horizontal-flip test-time augmentation, and
post-processed exactly as above.

\section{Evaluation}
\label{sec:eval}

\subsection{Protocol}
All local numbers use a local replication of the organiser's self-evaluation
notebook~\cite{notebook}: IoU$>0.5$ strictly, matched pairs counted, numerators
accumulating the IoU of the matched pairs, pooled over all
annotator--image rows; we also report the pooled pixel Dice.

The split is fixed for every experiment. Of the 707 training images, 296 carry
at least two annotators; 88 of them form the validation set (stratified by era
and by annotator count, shuffled with a fixed seed), and the remaining 619
images are the training set. The two sets are disjoint (0 shared file names).
The validation set contains 220 annotator--image rows, 1{,}514 ground-truth
masks and 21 annotators. The models are trained on the 619 images only; the 88
validation images are never trained on, but they are used for checkpoint
selection and for every hyper-parameter scan below, which we quantify in
Section~\ref{sec:limits}.

\subsection{Main results}
Table~\ref{tab:main} lists the progression of our own systems under one
protocol. The hand-crafted consensus detector reaches PQ 0.2201 and Dice
0.4861. Learning from consensus pseudo-labels gives 0.2696 / 0.6216 with a
512-pixel model; a 1024-pixel model with flip TTA gives 0.3152 / 0.6409.
Averaging four seeds adds $+0.0218$ PQ (0.3370 / 0.6512) and the
morphological closing a further $+0.0079$ (0.3449 / 0.6547), with 554
predicted instances over the 88 images. The 512-pixel row is an earlier and
much smaller model (0.48M parameters, 90 epochs, 14.5 minutes on an RTX 3070
laptop GPU) whose training log records 473 training and 85 validation images:
at that time 558 of the 707 training images had been downloaded, so it saw
473 images against the 619 used by the 1024-pixel models; the step to 1024
pixels therefore changes resolution, capacity, schedule and data size together
and is not an ablation of resolution alone.

\begin{table}[t]
\footnotesize
\setlength{\tabcolsep}{3.6pt}
\caption{Validation results on the fixed 88-image split, using the local
replication of the official metric (PQ pooled over annotator--image rows, Dice
pooled over pixels; TP/FP/FN in the official convention). TTA = horizontal-flip
test-time augmentation.}
\label{tab:main}
\begin{tabular}{lrrrrr}
\toprule
System & PQ & Dice & TP & FP & FN \\
\midrule
Hand-crafted detector & 0.2201 & 0.4861 & 455 & 698 & 1{,}059 \\
U-Net, 512 px & 0.2696 & 0.6216 & 603 & 636 & 911 \\
U-Net, 1024 px + TTA & 0.3152 & 0.6409 & 704 & 637 & 810 \\
4-seed probability average & 0.3370 & 0.6512 & 758 & 654 & 756 \\
\quad + $7\times7$ closing (submitted) & \textbf{0.3449} & \textbf{0.6547} & 779 & 632 & 735 \\
\bottomrule
\end{tabular}
\end{table}

\subsection{Ablations, including the negative ones}
\label{sec:ablation}
Table~\ref{tab:abl} summarises the variants we tested. Three of them are
negative results that we report because they shaped the final design.
Splitting an instance by its Left/Right class label costs $-0.0224$ PQ
(0.3142 to 0.2917): the split produces a narrower fragment that misses the
annotator boundary, and the metric punishes both the new false positive and
the new false negative. Replacing connected components by a geometric
watershed is much worse ($-0.2019$, 0.3142 to 0.1123): the watershed
over-segments the foreground into hundreds of fragments and the false-positive
count grows from 635 to 1{,}427. Widening the network from 32 to 48 base
channels also hurts: the base-48 run reaches 0.2096 against 0.2966 for the
base-32 run of the same family. The two runs differ in seed, so the comparison
also carries the seed spread quantified below, but this is the only wider
configuration we trained, and it ends 0.087 PQ behind while costing 2.25$\times$
the parameters and 48.7 instead of 29.0 minutes per seed, so we abandoned that
line.

Two further observations qualify the remaining gains. First, the closing kernel
is not a knife-edge: kernels 7, 9 and 13 give 0.3449, 0.3443 and 0.3441, a
spread of 0.0008, so the gain comes from closing the gaps rather than from the
particular kernel size. Second, the seed spread is much larger than the
ensembling gain: over the four seeds for which we have complete sweeps, the
best validation PQ ranges from 0.2070 to 0.2905 without flip TTA and from
0.2231 to 0.3142 with it, i.e. an envelope of 0.2070--0.3142, while the
four-seed average gains $+0.0218$ over the single-seed system it replaces.
Seed selection therefore carries a risk that ensembling only partly offsets.

\begin{table}[t]
\footnotesize
\setlength{\tabcolsep}{3.2pt}
\caption{Ablations on the same split, with the same probability maps and
evaluator; $\Delta$ (in brackets) is the change in PQ with respect to the row it
is compared against; TTA = horizontal-flip test-time augmentation,
$k$ = closing kernel.}
\label{tab:abl}
\begin{tabular}{lrr}
\toprule
Variant & PQ & TP/FP/FN \\
\midrule
TTA off (1 seed) & 0.2922 [$-0.0230$] & 689/772/825 \\
TTA on (1 seed) & 0.3152 & 704/637/810 \\
Closing $k=3/5/7/9/13$ & 0.3386--0.3449 & --- \\
Class split (Left/Right) & 0.2917 [$-0.0224$] & 649/669/865 \\
Watershed (best of 4) & 0.1123 [$-0.2019$] & 286/1{,}427/1{,}228 \\
Wider net (48 base ch.) & 0.2096 [$-0.0870$] & 494/942/1{,}020 \\
Best single seed & 0.3268 & --- \\
4-seed avg., no closing & 0.3370 [$+0.0102$] & 758/654/756 \\
4-seed avg. + closing & 0.3449 [$+0.0079$] & 779/632/735 \\
\bottomrule
\end{tabular}
\end{table}

\subsection{Public leaderboard}
Eight files have been submitted. Their public scores rise together with our own
local scores: 0.15 (first hand-crafted file, 126 of 180 test images covered),
0.16, 0.17 and 0.18 (the hand-crafted consensus detector, 174 of 180 images
covered), then 0.21 for the 512-pixel U-Net (local validation PQ 0.2696), 0.25
for a two-seed average (0.2993), 0.27 for the single 1024-pixel seed (0.3152)
and 0.29 for the submitted four-seed ensemble with closing (0.3449).
Dividing the public score by the local validation PQ of the same file gives
0.779, 0.835, 0.857 and 0.841 for
the four model submissions, i.e. a band of 0.78 to 0.86, which indicates that
the local protocol ranks our own variants in the same order as the public
board does.
We stress two caveats. The public score is a coarse signal, not an estimate of
our test performance: the API metadata announces a panoptic-quality
algorithm while the challenge page describes a mean-Dice score, and the
leaderboard exposes a single score column, so the two can never be reconciled
from the outside. And the local numbers behind those ratios are themselves
optimistic, for the reason given next.

\subsection{Reproducibility}
Every number quoted here is read from a retained evaluation artifact rather than
transcribed: the per-configuration evaluation files
(\fn{eval\_*\_val88.json}), the seed table (\fn{seed\_pq\_table.json},
15 training runs), the post-processing sweep
(\fn{postproc\_sweep\_val88.csv/.json}, 60 variants), the ensemble sweeps
(\fn{ens\_specs\_val88.csv}, \fn{custom\_specs\_val88.csv}, 64 and
28 configurations), the submission reports (\fn{submission\_*\_report.json}),
the captured public-leaderboard listing (\fn{kaggle\_submissions\_20261101.txt},
8 submissions) and the training log of the 512-pixel model
(\fn{train\_v1.log}). The file \fn{report\_evidence.json} maps the evaluation
numbers quoted here to their source artifact.

\subsection{Limitations}
\label{sec:limits}
Four limitations qualify everything above. \emph{(i) Selection optimism.} The
88 validation images play three roles: they monitor training and pick the
checkpoint, they select the threshold, the closing kernel and the area floor
over a 60-variant post-processing scan, and they select the seed. Every
validation number in this paper is therefore a post-selection maximum. We
quantified the effect by selecting the best variant on one half of the split
and evaluating it on the other half: the held-out mean is 0.3019 against
0.3172 for the best variant chosen on the full split, i.e. an optimism of
about 0.015 PQ, and the estimate ranges from $+0.002$ to $+0.015$ across four
different splits of the validation set. The same applies to the headline 0.3449, which was also
chosen on this split. \emph{(ii) Metric ambiguity.} The official description
of the leaderboard score is inconsistent between the API metadata and the
challenge page, so we use the public score only for ordering. \emph{(iii) Seed
variance.} The spread between the four seeds with complete sweeps
(0.2070--0.3142) is larger than most of
the individual improvements we report, and 6 of 15 training runs diverged with
a non-finite loss, which shows that the optimisation itself is fragile.
\emph{(iv) Data scale.} The supervision is limited to 707 public training
images; the 512-pixel model was trained on 473 of them, since only 558 were
on disk at the time and 85 of those were held out for validation, whereas the
1024-pixel models used all 619 images of the training split, so the effect of
the larger label set is not isolated.

\begin{acks}
We thank the organisers of the IEEE BigData Cup 2026 for the Solar Filament
Segmentation Challenge and for the public self-evaluation notebook that defines
the metric~\cite{notebook}. The data are the MLEcoFi 2024 dataset (version 1.0)
from the MLEcoFi Team (Earth-Space AI Research Lab with the National
Solar Observatory, Georgia State University), obtained through the Kaggle
competition files~\cite{competition}; the images originate from the
NSO/GONG network. We also thank the annotators whose multiple polygon
drawings per image make the consensus analysis in
Section~\ref{sec:consensus} possible.
\end{acks}

\bibliographystyle{ACM-Reference-Format}
\bibliography{main}

\end{document}
